\documentclass[11pt, aspectratio=169, table]{beamer}
  % papersize={16cm,9cm},

\mode<presentation>{}

%%%%%%%%%%%%%%%%%%%%%%%%%%%%%%
% load packages
% add packages if needed

%\usepackage[T1]{fontenc}
\usepackage[utf8]{inputenc}
\usepackage{xcolor}
\usepackage[english]{babel}
\usepackage{amsmath,amsfonts,amssymb,amsthm}
\usepackage{color}
\usepackage[tracking=smallcaps, letterspace=-55]{microtype} % package for font spacing
\usepackage[outline]{contour}
%\usepackage[perpage]{footmisc}
\usepackage{transparent} % for setting opacity of pictures
\usepackage{pifont} % for checkmark (\ding{51}) and cross (\ding{55})
\usepackage{booktabs,array} % for tables
\usepackage{graphicx} % for figures
\usepackage{tikz}
\usepackage{pgfpages}
\usepackage{ulem}
\usepackage{dirtytalk}
\usepackage{minted}
\usepackage{textcomp}
\usepackage{mwe}

\usemintedstyle{emacs}
\setminted{fontsize=\footnotesize,python3=true}
\definecolor{bg}{rgb}{0.97,0.97,0.95}
\setmintedinline{bgcolor=bg}
\setminted{bgcolor=bg}

%\graphicspath{} % to set the path of figures

\usetikzlibrary{shapes.multipart,positioning,matrix,external,shadows}

\renewcommand{\thefootnote}{\fnsymbol{footnote}}
\renewcommand{\thempfootnote}{\fnsymbol{mpfootnote}}

%\setbeameroption{show notes on second screen=bottom}

%%%%%%%%%%%%%%%%%%%%%%%%%%%%%%
% Fill in the following information: AUTHOR, TITLE, DATE

\author[Saska Dönges]{Saska Dönges}
\title[C++ Wut]{C++ basics}
\subtitle{}
\date{March 23. 2023, BSCS2011}%{\today}
\institute{Dept.\ of Computer science, University of Helsinki
}

%%%%%%%%%%%%%%%%%%%%%%%%%%%%%%
% set beamer colors

\definecolor{hyblue}{RGB}{0,155,255}
\setbeamercolor{alerted text}{fg=hyblue}
\setbeamercolor{structure}{fg=hyblue}

\setbeamertemplate{itemize item}{\color{hyblue}}

%%%%%%%%%%%%%%%%%%%%%%%%%%%%%%
% set font (helvetica plays the role of Arial)
\usepackage{helvet}
\renewcommand{\familydefault}{\sfdefault}

%%%%%%%%%%%%%%%%%%%%%%%%%%%%%%
% set frametitle

\setbeamercolor{frametitle}{fg=black}%{fg=hyblue}
\setbeamerfont{frametitle}{series=\bfseries, shape=\scshape, size=\huge}
\setbeamertemplate{frametitle}[default][left,leftskip=3.5cm] % left shift of frame title
\addtobeamertemplate{frametitle}{\vspace{0.5cm}}{\vspace{1cm}} % spacing above and below frame title

%%%%%%%%%%%%%%%%%%%%%%%%%%%%%%
% set footline and headline
\beamertemplatenavigationsymbolsempty
\setbeamertemplate{headline}{ }
\setbeamertemplate{footline}{%
	 \usebeamercolor[fg]{page number in head/foot}%
	 \usebeamerfont{page number in head/foot}%
	\hspace{0.5cm}	
	\includegraphics[width=2.5cm]{HY__LC05_txt__L_3L_B3____BW_cropped}
	\hfill
	\insertshorttitle\ /\ \insertshortauthor	\hfill
	\insertdate	\hspace{0.5cm}
	\insertframenumber\,/\,\inserttotalframenumber \hspace{0.5cm} \vskip2pt%
}

%%%%%%%%%%%%%%%%%%%%%%%%%%%%%%
% Other style settings
\useinnertheme{circles}

%%%%%%%%%%%%%%%%%%%%%%%%%%%%%%
% For video
\usepackage{media9}%
\newcommand{\includemovie}[3]{%
\includemedia[%
width=#1,height=#2,%
activate=onclick,%
deactivate=pageclose,%
addresource=#3,%
flashvars={%
src=#3 % same path as in addresource!
%&autoPlay=true % default: false; if =true, automatically starts playback after activation (see option ‘activation)’
&loop=true % if loop=true, media is played in a loop
&controlBarAutoHideTimeout=0 %  time span before auto-hide
}%
]{}{StrobeMediaPlayback.swf}%
}

%%%%%%%%%%%%%%%%%%%%%%%%%%%%%%
%%%%%%%%%%%%%%%%%%%%%%%%%%%%%%

\begin{document}

%% Title slide with HY picture
%{
%\usebackgroundtemplate{
%\setlength{\unitlength}{1cm}
%\begin{picture}(16,9)
%\transparent{0.33}
%\put(0.3,0.8){\includegraphics[width=15cm,height=7.8cm]{hy_brand_image_1_wide_background_hires_logo} }
%\end{picture}
%}
%\setbeamertemplate{headline}{ }%
%\begin{frame}{}
%\vspace{2.5cm}
%\begin{center}
%\huge
%{\bf \inserttitle } \\
%{\large \insertauthor}
%\end{center}
%\end{frame}
%}



% Title page without picture
{
\usebackgroundtemplate{
\setlength{\unitlength}{1cm}
\begin{picture}(16,9)
	\put(-0.1,5){ \includegraphics[width=4cm]{HY__LA01_Flame_____B3____BW} }
\end{picture}
\setlength{\unitlength}{1pt}
}
\setbeamertemplate{headline}{ }%
\begin{frame}[noframenumbering]{}
\vspace{2.5cm}
\begin{center}
\huge
\textcolor{hyblue}{ \bf
\inserttitle
} \\
{\large
\insertauthor}
\end{center}
\end{frame}
}
%

%%%%%%%%%%%%%%%%%%%%%%%%%%%%%%
%% Main text

%% set HY Logo on the upper-left corner:
\usebackgroundtemplate{
\setlength{\unitlength}{1cm}
\begin{picture}(16,2.5)
\put(0.1,0){
\includegraphics[width=2.5cm]{HY__LA01_Flame_____B3____BW}\hfill
}
\end{picture}
\setlength{\unitlength}{1pt}
}
%%

%%%%%%%%%%%%%%%%%%%%%%%%%%%%%%
%% Modify text below

\begin{frame}{Agenda}
\setlength\parskip{\fill}
\tableofcontents
\alert{Interrupt} to comment or ask questions! You can get \alert{points}!
\end{frame}

\section{On a high level}
\begin{frame}{What is C++}
\setlength\parskip\fill
C++ is a programming language created in 1979, as an extension to C.

\pause
\only<2>{C++ is a procedural, imperative, functional, object-oriented, generic, modular programming language with support for metaprogramming.}

\pause
\sout{C++ is a procedural, imperative, functional, object-oriented, generic, modular programming language with support for metaprogramming.}

\only<3>{C++ is a federation of multiple concepts. C, C with classes, Template C++ and, STL... Among others.}

\pause
\sout{C++ is a procedural, object-oriented, functional, generic programming language with support for metaprogramming.}

C++ is a cluster****. 20 different languages in a trenchcoat pretending to be one well-adjusted language.

\pause
And probably the most used systems programming language in the world today.
\end{frame}

\begin{frame}{The (main) parts (you should care about)}
\setlength\parskip\fill
Any valid C code is valid C++ code.

C with classes. Instead of just structs in C, we actually have classes with member functions. YAY.

C++ templating. Compile-time generics. Introduces the idea of doing computation at compile time. 
And the neat but mindboggling art of template metaprogramming. Don't ask me about template metaprogramming.

STL (standard template library). A full complement of \sout{fast} not too slow data structures and algorithms provided 
as standard.
\end{frame}

\section{Hello world!}
\begin{frame}[fragile]{Hello world!}
\begin{minted}{cpp}
#include <stdio.h>
#include <iostream>

int main(int argc, char const *argv[]) {
    std::string hello = "Hello, world!\n";
    puts("Hello, world!");
    std::fputs(hello.c_str(), stdout);
    printf(hello.c_str());
    std::cout << hello << std::flush;
    std::cout.write(hello.c_str(), hello.size());
    return 0;
}
\end{minted}
\end{frame}

\begin{frame}[fragile]{Hello world!}
Or simply

\begin{minted}{cpp}
#include <iostream>

int main() {
    std::cout << "Hello, world!" << std::endl;
}
\end{minted}

Why don't we even get a warning for not returning enything?
\end{frame}

\begin{frame}[fragile]{Hello world!}
The main arguments are possibly familiar

\begin{minted}{cpp}
int main(int argc, char const *argv[]) {
	std::cout << "Hello, " << argv[1] << "!" << std::endl;
}
\end{minted}

\begin{minted}{shell}
~/scratch$ ./hello Exactum
Hello, Exactum!
~/scratch& ./hello
\end{minted}
\pause\vspace{-1.2cm}
\begin{minted}{shell}
Hello, ~/scratch$
\end{minted}

Or \texttt{Segmentation fault}, depending on stuff and things.
\end{frame}

\section{Undefined behaviour}
\begin{frame}[fragile]{Array index out of bounds}
\begin{minted}{cpp}
int main(int argc, char const *argv[]) {
    int bla[10];
    for (int i = 0; i < 100; i++) {
        std::cout << bla[i]-- << " " << std::flush;
    }
}
\end{minted}

\pause\begin{minted}{shell}
-1226429632 32689 0 0 -872412576 21939 -872413072 
21939 82839072 32765 -2017372928 2006707886 -872412576 
21939  ... 43551 32765 82843582 32765 82843595 
32765 *** stack smashing detected ***: <unknown> terminated
Aborted
\end{minted}
\end{frame}

\begin{frame}[fragile]{Array index out of bounds}
\begin{minted}{cpp}
int main(int argc, char const *argv[]) {
    int* bla = new int[10]();
    for (int i = 0; i < 100; i++) {
        std::cout << ++bla[i] << " " << std::flush;
    }
}
\end{minted}

\pause\begin{minted}{shell}
1 1 1 1 1 1 1 1 1 1 1042 1 842276914 825833016 8247 1 1 1 1 1 
1 1 1 1 1 1 1 1 1 1 1 ... 1 1 1 1 1 1 1 1 1 1 1 1 1 1 1 1 1 1 
\end{minted}

Or \texttt{Segmentation fault} if the OS feels like it.
\end{frame}

\begin{frame}{Array index out of bounds}
\setlength\parskip{\fill}
Don't assume any bounds checking is done in C++. 

\texttt{stl} containers typically do bounds checking for \texttt{at} 
functions but not for angle bracket access. \mintinline{cpp}{std::vector.at()} is safer but slower.

What happens when you attempt to access memory outside of where you should is undefined.

If the memory is "owned" by you, and you are not smashing the stack the program will not crash even
if everything is completely broken.

If you attempt to access memory you are not allowed to, the OS will segfault the program.
\end{frame}

\begin{frame}[fragile]{Valgrind to the rescue}

For the second version:

\begin{minted}{shell}
~/scratch$ valgrind ./hello
==2799== Command: ./hello
==2799==
1 1 1 1 1 1 1 1 1 1 ==2799== Invalid read of size 4
==2799==    at 0x10898F: main (in /home/saskeli/scratch/hello)
\end{minted}

Valgrind may not detect stack smashing.
\end{frame}

\begin{frame}[fragile]{More undefined behaviour}
\setlength\parskip{2pt}
Undefined behaviour is any time something happens for which the C++ standard has a correct behaviour.

A typical example is reading memory that is not allocated to the process. Reading and writing may or 
may not work as expected, and the process may or may not mangle itself or get killed by the OS.

Most of the time undefined behaviour encountered in the wild is due to memory errors. 

Another trap for new player is the bit shift where 

\vspace{-0.3cm}\begin{minted}{cpp}
uint32_t x = 1;
x = x << n;
\end{minted}

\vspace{-0.5cm}for any $n \geq 32$ is undefined behaviour.
\end{frame}

\section{Memory management}
\begin{frame}[fragile]{Memory management}
Recall

\begin{minted}[highlightlines={3}]{cpp}
int main(int argc, char const *argv[]) {
    int* bla = new int[10]();
    for (int i = 0; i < 10; i++) {
        std::cout << ++bla[i] << " " << std::flush;
    }
}
\end{minted}

but now fixed
\end{frame}

\begin{frame}[fragile]{Memory management}
\vspace{-0.5cm}
\begin{minted}[highlightlines={5,9}]{shell}
==2863== Command: ./hello
==2863==
1 1 1 1 1 1 1 1 1 1 ==2863==
==2863== HEAP SUMMARY:
==2863==     in use at exit: 40 bytes in 1 blocks
==2863==   total heap usage: 3 allocs, 2 frees, 73,768 bytes allocated
==2863==
==2863== LEAK SUMMARY:
==2863==    definitely lost: 40 bytes in 1 blocks
==2863==    indirectly lost: 0 bytes in 0 blocks
==2863==      possibly lost: 0 bytes in 0 blocks
==2863==    still reachable: 0 bytes in 0 blocks
==2863==         suppressed: 0 bytes in 0 blocks
==2863== Rerun with --leak-check=full to see details of leaked memory
\end{minted}
\end{frame}

\begin{frame}[fragile]{Memory management}
Not really important here, since program termination \linebreak
will free up to leaked memory but...

\begin{minted}[highlightlines={6}]{cpp}
int main(int argc, char const *argv[]) {
    int* bla = new int[10]();
    for (int i = 0; i < 10; i++) {
        std::cout << ++bla[i] << " " << std::flush;
    }
    delete(bla);
}
\end{minted}

Every \texttt{new} should have a matching \texttt{delete}, \linebreak
and every \texttt{malloc}/\texttt{calloc} a matching \texttt{free}.
\end{frame}

\begin{frame}[fragile]{Memory leaks}
Especially important when you have classes that allocate memory internally. For example, we could 
have a neat dynamic array that keeps track of the sum of contained values:

\begin{minted}{cpp}
template<class dtype>
class SumVector {
   private:
    uint64_t capacity;
    uint64_t size;
    uint64_t sum;
    dtype* data;
\end{minted}
\end{frame}

\begin{frame}[fragile]{Memory leaks (interlude)}
Actually you'd do the following unless you have a very good reason not to.

\begin{minted}{cpp}
template<class dtype>
class SumVector {
   private:
    uint64_t sum;
    std::vector<dtype> data;
\end{minted}

And after this the compiler can probably take care of all memory management automatically.
\end{frame}

\begin{frame}{Memory leaks}
\setlength\parskip\fill
It's common for objects to be (de)allocated ``under the hood'' when doing assignments and passing objects around.

Any time this happens, there is a risk deallocation of an object with reference members leaves orphaned memory 
that is \alert{permanently lost} and only deallocated when the process exits.

This leads us conveniently to the next part
\end{frame}

\section{Copy and move semantics}
\begin{frame}[fragile]{Copy and move semantics (soon)}
When is memmory allocated\footnote{Perhaps not actually, but for practical purposes} for member variables and storage for \texttt{vec}?

\begin{minted}{cpp}
int something_with_a_vector(T& param) {
    /* loads of omitted code here */
    std::vector<int> vec;
    /* loads of code using vec here */
    return x;
}
\end{minted}
\end{frame}

\begin{frame}[fragile]{Copy and move semantics (soon)}
Memory for local variables is reserved on stack as part of the memory footprint of the function. 
This includes the local member variables of the vector.

\begin{minted}{cpp}
int something_with_a_vector(T& param) {
    /* loads of omitted code here */
    std::vector<int> vec; // Storage space allocated here (or earlier)
    /* loads of code using vec here */
    return x;
} //Storage space is actually (typically) deallocated here.
\end{minted}

Mhen vec goes out of scope the member variables will be outside the stack 
but the data will still be there. \alert{VERY DANGEROUS!}
\end{frame}

\begin{frame}[fragile]{Copies}
Any time you assign a value to a new variable or pass by value to a function, a copy will be made.

\begin{minted}{cpp}
T thing; // Some heavy object
T other = thing // Copy of thing

void some_function(T thing) // Will create a copy of what is passed in.
\end{minted}

This can be potentially costly and dangerous.
\end{frame}

\begin{frame}[fragile]{Copies}
\setlength\parskip\fill
Any time a copy of a class \texttt{T} is created, the copy constructor 
\mintinline{cpp}{T(const T& other)} will be called.

If you have not created one or forbidden the creation of one \alert{one will be provided by the comipiler}.

The default copy constructor \alert{only copies members} naively.
\end{frame}

\begin{frame}[fragile]{Copies}
So for our neat \texttt{SumVector}:

\begin{minted}{cpp}
int main() {
    SumVector first;
    first.push_back(1).push_back(2).push_back(3);
    SumVector second = first;
    second.pop().push_back(66).push_back(102);
    first.print();
    second.print();
}
\end{minted}

What happens?
\end{frame}

\begin{frame}[fragile]{Broken copies}
\begin{minted}{shell}
~/scratch$ ./hello
Size = 3
Sum = 6
1 2 66
Size = 4
Sum = 171
1 2 66 102
free(): double free detected in tcache 2
Aborted
~/scratch$
\end{minted}
\end{frame}

\begin{frame}[fragile]{Fixed copies}
To make copies work correctly we need to manually deal with the allocated \texttt{data}.

\begin{minted}{cpp}
SumVector(const SumVector& other)
    : capacity(other.capacity), size(other.size), sum(other.sum) {
    data = (int*)malloc(capacity * sizeof(int));
    std::memcpy(data, other.data, size * sizeof(int));
}
\end{minted}

Or better yet \mintinline{cpp}{SumVector(const SumVector&) = delet;}, so the compiler stops any copying.
\end{frame}

%{
%\usebackgroundtemplate{
%\transparent{0.3} % requires \usepackage{transparent} & compile twice
%\setlength{\unitlength}{1cm}
%\begin{picture}(16,9)
%\put(0.2,1){
%\includegraphics[width=15cm,height=7.5cm]{hy_brand_image_1_wide_background_hires_logo}
%}
%\end{picture}
%\setlength{\unitlength}{1pt}
%}
%
%\begin{frame}{\vspace{3cm}That's it}
%\begin{center}
%\vspace{-4cm}{\huge Questions?}
%\end{center}
%\end{frame}
%}

\end{document}
